\section*{Hoofdstuk \ref{ch:g11:dissolution} --- Ionaire en moleculaire vaste stoffen oplossen}

\begin{solution}{exo:g11:dissolution:1}
\ce{NaCl(s) -> Na+(aq) + Cl-(aq)};
\ce{CaCl2(s) -> Ca^{2+}(aq) + 2Cl-(aq)};
\ce{Na2SO4(s) -> 2Na+(aq) + SO4^{2-}(aq)};
\ce{AlCl3(s) -> Al^{3+}(aq) + 3Cl-(aq)}.
\end{solution}

\begin{solution}{exo:g11:dissolution:2}
$[\ce{Na+}] = 2 \times 0.050 = \qty{0.10}{mol/L}$;
$[\ce{SO4^{2-}}] = \qty{0.050}{mol/L}$.
\end{solution}

\begin{solution}{exo:g11:dissolution:3}
\omterm{def:g11:dissolution:solvation}{Hydratatie} is het omringd worden van het \omterm{def:g9:ions:ion}{ion} door watermoleculen die eraan
vastgehouden worden. Het zuurstofuiteinde ($\delta^-$) wijst naar een
\omterm{def:g9:ions:ion}{kation}, het waterstofuiteinde ($\delta^+$) naar een \omterm{def:g9:ions:ion}{anion}: tegengestelde
ladingen trekken elkaar aan.
\end{solution}

\begin{solution}{exo:g11:dissolution:4}
$M(\ce{AlCl3}) = 27.0 + 3 \times 35.5 = \qty{133.5}{g/mol}$;
$n = 2.67 / 133.5 = \qty{0.0200}{mol}$; $c = 0.0200 / 0.2000 =
\qty{0.100}{mol/L}$; $[\ce{Al^{3+}}] = \qty{0.100}{mol/L}$,
$[\ce{Cl-}] = \qty{0.300}{mol/L}$.
\end{solution}

\begin{solution}{exo:g11:dissolution:5}
De kop \ce{-COO-} is ionair en wordt door water aangetrokken: \omterm{def:g11:dissolution:hydrophilic}{hydrofiel}.
De staart van 17 koolstofatomen, met alleen \ce{C-C}- en bijna apolaire
\ce{C-H}-bindingen, is apolair: \omterm{def:g11:dissolution:hydrophilic}{hydrofoob}.
\end{solution}

\begin{solution}{exo:g11:dissolution:6}
Zout: water (ionaire vaste stof, polair \omterm{def:g6:solutions-and-solubility:solute}{oplosmiddel}). Kaarsvet:
cyclohexaan (apolaire ketens). Suiker: water (zijn \ce{O-H}-groepen vormen
\omterm{def:g11:polarity-and-cohesion:hydrogen-bond}{waterstofbruggen} met water). Jood: cyclohexaan (\omterm{def:g11:polarity-and-cohesion:polar-molecule}{apolair molecuul}).
\end{solution}

\begin{solution}{exo:g11:dissolution:7}
De \ce{O-H}-groep van ethanol vormt \omterm{def:g11:polarity-and-cohesion:hydrogen-bond}{waterstofbruggen} met watermoleculen,
die de bruggen vervangen die ze verbreekt: ethanol past in water. Hexaan,
apolair, kan zulke bindingen niet vormen; de watermoleculen, onderling
verbonden door \omterm{def:g11:polarity-and-cohesion:hydrogen-bond}{waterstofbruggen}, laten het niet toe.
\end{solution}

\begin{solution}{exo:g11:dissolution:8}
De \omterm{def:g11:dissolution:hydrophilic}{hydrofobe} staarten mijden het water en \omterm{def:g2:mixing-and-dissolving:dissolve}{lossen op} in het vet, in het
midden; de geladen koppen worden door water aangetrokken en blijven
buiten. Alle \omterm{def:g11:dissolution:amphiphilic}{micellen} dragen negatieve ladingen aan hun oppervlak, en
gelijke ladingen stoten elkaar af: ze blijven uit elkaar.
\end{solution}

\begin{solution}{exo:g11:dissolution:9}
Cyclohexaan: het is niet \omterm{def:g6:solutions-and-solubility:miscible}{mengbaar} met water, dus het vormt een aparte laag
die je kunt aftappen, en de geurstof (apolair) lost er goed in op. Ethanol
is \omterm{def:g6:solutions-and-solubility:miscible}{mengbaar} met water: er ontstaat geen tweede laag, en er valt niets te
scheiden.
\end{solution}

\begin{solution}{exo:g11:dissolution:10}
Cyclohexaan is minder dicht dan water. De onderste, waterige laag wordt
via de kraan afgetapt. Het jood zit in de bovenste laag, het cyclohexaan,
die in de trechter blijft.
\end{solution}

\begin{solution}{exo:g11:dissolution:11}
Niets: kopersulfaat is ionair en lost niet op in het apolaire
cyclohexaan. Het water blijft blauw en het cyclohexaan kleurloos.
\end{solution}

\begin{solution}{exo:g11:dissolution:12}
$[\ce{Cl-}] = 2c$, dus $c = \qty{0.150}{mol/L}$;
$n = 0.150 \times 0.2500 = \qty{0.0375}{mol}$;
$m = 0.0375 \times 111.1 = \qty{4.17}{g}$.
\end{solution}

\begin{solution}{exo:g11:dissolution:13}
De calcium- en magnesiumionen slaan de stearaationen neer als kalkzeep:
de zeep is op voordat ze \omterm{def:g11:dissolution:amphiphilic}{micellen} kan vormen, en schuimt slecht. Alleen
calcium: \qty{0.010}{mol} \ce{Ca^{2+}} neemt \qty{0.020}{mol} stearaat in
beslag, dat is $0.020 \times 306.0 =
\qty{6.1}{g}$ natriumstearaat per liter.
\end{solution}

\begin{solution}{exo:g11:dissolution:14}
Volume \qty{0.2000}{L}. \ce{Na+}: $0.020 / 0.2000 = \qty{0.10}{mol/L}$;
\ce{Ca^{2+}}: $0.010 / 0.2000 = \qty{0.050}{mol/L}$; \ce{Cl-}:
$(0.020 + 0.020) / 0.2000 = \qty{0.20}{mol/L}$. Positieve lading:
$0.10 + 2 \times 0.050 = 0.20$; negatieve: $0.20$. Neutraal.
\end{solution}

\begin{solution}{exo:g11:dissolution:15}
Een liter heptaan heeft een massa van \qty{0.68}{kg} en kan
$17 \times 0.68 = \qty{12}{g}$ jood bevatten, ongeveer 40 keer meer dan
een liter water. Als je ze samen schudt, gaat bijna al het jood over naar
het heptaan: de \omterm{def:g10:chemical-species:extraction}{extractie} is doeltreffend.
\end{solution}

\begin{solution}{pb:g11:dissolution:1}
\textbf{1.} Uit de gesteenten waar het water doorheen is gestroomd
(kalksteen, dolomiet, gips), die een beetje \omterm{def:g2:mixing-and-dissolving:dissolve}{oplossen}.

\textbf{2.} $[\ce{Ca^{2+}}] = 0.120 / 40.1 = \qty{2.99e-3}{mol/L}$;
$[\ce{Mg^{2+}}] = 0.024 / 24.3 = \qty{9.9e-4}{mol/L}$.

\textbf{3.} \ce{Ca(HCO3)2 -> Ca^{2+}(aq) + 2HCO3-(aq)}; dus
$[\ce{HCO3-}] = 2 \times \num{2.99e-3} = \qty{5.99e-3}{mol/L}$.

\textbf{4.} $\num{2.99e-3} \times 150 = \qty{0.449}{mol}$.

\textbf{5.} $18 \times 12.0 + 35 \times 1.0 + 2 \times 16.0 + 23.0 =
\qty{306.0}{g/mol}$.

\textbf{6.} \ce{C17H35COONa(s) -> C17H35COO-(aq) + Na+(aq)}.

\textbf{7.} Kop \ce{-COO-}: \omterm{def:g11:dissolution:hydrophilic}{hydrofiel}; staart \ce{C17H35-}: \omterm{def:g11:dissolution:hydrophilic}{hydrofoob}.
Het heeft beide soorten delen: \omterm{def:g11:dissolution:amphiphilic}{amfifiel}.

\textbf{8.} Een piepklein bolletje van zeepionen, staarten naar binnen,
koppen naar buiten. De staarten \omterm{def:g2:mixing-and-dissolving:dissolve}{lossen op} in het vet, dat uiteenvalt in
druppeltjes die in \omterm{def:g11:dissolution:amphiphilic}{micellen} ingepakt zijn; de geladen oppervlakken houden
de druppeltjes in het water uit elkaar, en het water spoelt ze weg.

\textbf{9.} De zeep moet \omterm{def:g9:ions:ion}{ionen} en \omterm{def:g11:dissolution:amphiphilic}{micellen} vormen, en daarvoor is water
rond de koppen nodig; en het vet wordt door het spoelwater afgevoerd. In
olie zou het vet gewoon opgelost op de huid blijven.

\textbf{10.} \ce{2C17H35COO- + Ca^{2+} -> Ca(C17H35COO)2}: ladingen
$2 \times (-1) + 2 = 0$ links, $0$ rechts.

\textbf{11.} \ce{C17H35COO-}: overmaat $- 2x$; \ce{Ca^{2+}}:
$\num{2.99e-3} - x$; calciumstearaat: $x$. De calciumionen zijn
beperkend: $x_{\max} = \qty{2.99e-3}{mol}$.

\textbf{12.} $2 \times \num{2.99e-3} = \qty{5.99e-3}{mol}$ per liter.

\textbf{13.} $M = 36 \times 12.0 + 70 \times 1.0 + 4 \times 16.0 + 40.1 =
\qty{606.1}{g/mol}$; $\num{2.99e-3} \times 606.1 = \qty{1.81}{g}$
kalkzeep per liter.

\textbf{14.} $2 \times \num{9.9e-4} = \qty{1.98e-3}{mol}$ per liter.

\textbf{15.} $2 \times 0.449 = \qty{0.898}{mol}$.

\textbf{16.} $0.898 \times 306.0 = \qty{275}{g}$.

\textbf{17.} Magnesium: $\num{9.9e-4} \times 150 = \qty{0.148}{mol}$, dat
\qty{0.296}{mol} stearaat in beslag neemt, \qty{91}{g} zeep; in totaal
ongeveer $275 + 91 = \qty{366}{g}$.

\textbf{18.} $275 / 100 \approx 2.7$ stukken.

\textbf{19.} Het calcium van één bad laat ongeveer \qty{0.27}{kg} zeep
verloren gaan als kalkzeep.
\end{solution}
